\section{Sections, divided powers, and the envelope of \S50}
\label{sec:E}

We use \(\operatorname{Env}(\Gamma)\) for the inverse system of \S50.4,
and \(E_\Gamma\) for its affine inverse limit.  We use
\(\underline\Gamma_M\) for the single finite type group scheme obtained
from a chosen normal lattice \(M\subset\Gamma\).  These are different
objects.  The distinction is already explicit in the manuscript:
\S50.3, fols.~656--658, C4 scans 17--21; \S50.4, fols.~659--662,
C4 scans 23--29; and \S51.1, fol.~665, C4 scan 35.

\subsection{What ``unipotent'' means in the question}

\begin{statement}{E.1}{Proved}{Two definitions that must not be identified}
In C4 scan 5, fol.~650bis, Grothendieck calls a geometric section
\(g\in G(\overline k)\) unipotent when
\(g=\phi(1)\) for a homomorphism
\(\phi:\mathbb G_{a,\overline k}\to G_{\overline k}\).
In characteristic \(p>0\), every such element satisfies \(g^p=1\).
Consequently, an element of order \(p^2\) in a unipotent algebraic group
is a counterexample to the assertion that all conventionally unipotent
elements arise this way, but is not a counterexample satisfying the
manuscript's own definition.  The quasi-unipotent condition in that
definition is equivalent, on geometric fibres of positive characteristic,
to finite order.
\end{statement}
\begin{proof}
The multiplication-by-\(p\) map on \(\mathbb G_a\) is zero in
characteristic \(p\), so \(\phi(1)^p=\phi(p)=1\).
If a power \(g^n\) lies in such an image, then \(g^{np}=1\).
Conversely, if \(g\) has finite order, a power is the identity and is
the value of the zero homomorphism.  This proves all assertions without
assuming that the group itself is connected, smooth, or affine.
\end{proof}

\begin{statement}{E.2}{Proved}{The characteristic-zero assertion for linear targets}
Let \(A\) be a \(\mathbb Q\)-algebra and
\(G\hookrightarrow\operatorname{GL}_{r,A}\) a closed immersion of affine
group schemes.  Suppose that \(g\in G(A)\) is fibrewise unipotent in the
sense of E.1.  There is a unique homomorphism of \(A\)-group schemes
\(\phi:\mathbb G_{a,A}\to G\) such that \(\phi(1)=g\).  In the specified
representation,
\[
 N=g-I,\qquad X=\log(I+N),\qquad \phi(t)=\exp(tX),
\]
where the two expressions are finite polynomials.  No reducedness
assumption on \(A\) is needed.
\end{statement}
\begin{proof}
On every geometric fibre, \(N\) is nilpotent.  The non-leading
coefficients of its characteristic polynomial therefore belong to the
nilradical of \(A\).  They generate a nilpotent ideal \(J\): an ideal
generated by finitely many nilpotent elements is nilpotent, by expanding
a sufficiently long product.  Cayley--Hamilton gives
\(N^r\in J\operatorname{End}(A^r)\), hence \(N^{rq}=0\) if \(J^q=0\).
The finite formal identities for logarithm and exponential now apply in
the nilpotent algebra generated by \(N\), and give a homomorphism to
\(\operatorname{GL}_{r,A}\) with value \(g\).

For every integer \(n\), this homomorphism takes the value \(g^n\in G(A)\).
If \(F\) belongs to the defining ideal of \(G\), its pullback is a
polynomial \(P(t)\in A[t]\) vanishing at all integers.  For a degree
bound \(d\), evaluation at \(0,\ldots,d\) has Vandermonde determinant
\(\prod_{i<j}(j-i)\), a unit in \(A\).  Thus \(P=0\), so the map factors
through \(G\).

Finally, a polynomial matrix homomorphism \(P:\mathbb G_a\to
\operatorname{GL}_r\) obeys
\(P'(t)=P(t)D\), where \(D=P'(0)\).  Comparing coefficients over a
\(\mathbb Q\)-algebra gives \(P(t)=\sum t^iD^i/i!\); polynomiality
implies that \(D\) is nilpotent.  Its value at 1 determines
\(D=\log P(1)\).  This proves uniqueness.  In particular, this proof
justifies the descent argument of C4 scan 6 for linear targets.
\end{proof}

\begin{statement}{E.3}{Proved}{The exact integral divided-power criterion}
Let \(A\) be a torsion-free \(\mathbb Z\)-algebra,
\(G\hookrightarrow\operatorname{GL}_{r,A}\) a closed subgroup scheme,
and \(g\in G(A)\) such that \(g-I\) is nilpotent.  Put
\(X=\log g\in M_r(A\otimes\mathbb Q)\).  A homomorphism
\(\mathbb G_{a,A}\to G\) taking 1 to \(g\) exists if and only if
\[
                  X^i/i!\in M_r(A)\quad\hbox{for every }i\geq 0.
\]
When it exists it is unique.  Thus the coefficient criterion (10) of
\S50.1, fol.~652, C4 scan 9, is correct, but the inference to it from
(9) is false.
\end{statement}
\begin{proof}
After tensoring with \(\mathbb Q\), E.2 forces the polynomial
\(\exp(tX)\).  Polynomial coefficients over the torsion-free ring \(A\)
are uniquely determined in \(A\otimes\mathbb Q\).  This proves
necessity and uniqueness.  If all coefficients belong to \(A\), the
exponential defines a polynomial matrix with inverse \(\exp(-tX)\);
the addition identity is a polynomial identity over \(A\), because it
holds after tensoring with \(\mathbb Q\).  The same injectivity applied
to each defining equation of \(G\) proves factorization through \(G\).
\end{proof}

\begin{statement}{E.4}{Refuted}{The actual exponential integrality implication fails in rank three}
For the nilpotent Jordan matrix \(J=E_{12}+E_{23}\in M_3(\mathbb Z)\), set
\[
 X=J-\tfrac12J^2
   =\begin{pmatrix}0&1&-1/2\\0&0&1\\0&0&0\end{pmatrix}.
\]
Then \(X^3=0\), \(\exp X=I+J\), and
\(\exp(-X)=I-J+J^2\) are integral, whereas \(X\) is not integral at 2.
The assertion in C4 scan 9 that the \(n=3\) case of (9) implies (10)
is false.  The strongest exact replacement is E.3, and E.10 below
proves the weaker assertion after taking a positive power.
\end{statement}
\begin{proof}
Here \(X^2=J^2\), so the two exponential sums are exactly the displayed
matrices.  The \((1,3)\)-entry of \(X\) is \(-1/2\).  In rank two,
\(X^2=0\) and \(\exp X=I+X\), so this particular implication cannot
fail in smaller rank.  The exact rational identities and both inverse
matrices are checked in \texttt{code/E\_integrality.py} and recorded in
\texttt{code/E\_integrality.json}.  Notice that \(I+J\) modulo 2 has
order 4.  By E.1 this example alone does not answer the stronger
fibrewise question; E.6 does.
\end{proof}

\begin{statement}{E.5}{Refuted}{A counterexample to the literal signed reading of (11)}
The literal implication (11)\(\Rightarrow\)(12) in C4 scans 9--10 is
false for \(X\in M_4(\mathbb Q)\), even over \(A=\mathbb Z_{(2)}\)
(\S50.1, fols.~652--652').
One counterexample is
\[
 X=\begin{pmatrix}0&1/2&0&0\\0&0&12&2\\0&0&0&1/2\\0&0&0&0\end{pmatrix},
 \quad
 A_2=\tfrac12X^2=3E_{13}+3E_{24}+\tfrac12E_{14},
 \quad A_3=\tfrac16X^3=\tfrac12E_{14}.
\]
Thus \(A_2+A_3\) and \(A_2-A_3\) are integral, but neither \(A_2\)
nor \(A_3\) is.  If one adds the missing hypothesis \(X\in M_4(A)\),
for a torsion-free \(\mathbb Z\)-algebra \(A\), the literal implication
is true.
\end{statement}
\begin{proof}
Direct multiplication gives the displayed formulas and \(X^4=0\).
The matrix \(X_0=2X\) is integral and not divisible by 2; hence this
also gives precisely the form requested in (11bis), with \(\nu=1\).
It was the first counterexample after 168 candidates in the explicitly
specified finite rational grid in \texttt{E\_integrality.py}.

For the correction, write \(C=X^2/2+X^3/6\).  Over \(A[1/2]\),
adding and subtracting the two signed expressions isolates both terms.
Over \(A[1/3]\), the matrix \(3I+X\) is invertible with inverse the
finite geometric series in \(-X/3\), and
\[
                    X^2/2=3C(3I+X)^{-1}
\]
is integral.  Subtracting from \(C\) gives \(X^3/6\).  The two open
sets cover \(\operatorname{Spec}A\), because \((2,3)=A\); membership
in \(M_4(A)\) is local, so the result follows.

This reading must be kept distinct from E.4.  The actual hypothesis
\(\exp(\pm X)\in M_4(A)\) contains \(\pm X\), whereas (11) omits it.
Visual inspection of scan 9 confirms that the omission is in the
manuscript, not a silent correction supplied here.
\end{proof}

\begin{statement}{E.6}{Refuted}{A genuinely fibrewise-unipotent integral section need not extend}
Let \(N\in M_5(\mathbb Z)\) have superdiagonal \((1,2,3,2)\) and all
other entries zero.  The section \(g=I+N\in\operatorname{GL}_5(\mathbb Z)\)
is fibrewise unipotent in the precise sense of E.1, but is not
\(\phi(1)\) for any \(\phi:\mathbb G_{a,\mathbb Z}\to
\operatorname{GL}_{5,\mathbb Z}\).  More generally, over
\(\mathbb Z_{(p)}\), the implication
\[
 N^p\in pM_r(\mathbb Z_{(p)}),\qquad N^r=0
 \quad\Longrightarrow\quad
 (I+N)^t\in M_r(\mathbb Z_{(p)}[t])
\]
holds for \(r\leq p^2\), and the bound is optimal uniformly in \(N\).
The first counterexample from the weighted-shift family occurs in
degree \(p^2\), or rank \(p^2+1\).
\end{statement}
\begin{proof}
For the rank-five matrix,
\(N^2\equiv0\pmod2\), \(N^3\equiv0\pmod3\), and \(N^5=0\).
Over a field of characteristic \(p\), a matrix \(I+\overline N\) with
\(\overline N^p=0\) lies on a one-parameter subgroup: the truncated
logarithm \(D=\sum_{j=1}^{p-1}(-1)^{j+1}\overline N^j/j\) satisfies
\(D^p=0\), and the truncated exponential \(\sum_{j=0}^{p-1}t^jD^j/j!\)
has the required group law and value.  In characteristic zero use E.2.
This treats every geometric fibre of the integral section.

Nevertheless,
\[
 X=\log(I+N)=
 \begin{pmatrix}
 0&1&-1&2&-3\\0&0&2&-3&4\\0&0&0&3&-3\\0&0&0&0&2\\0&0&0&0&0
 \end{pmatrix},
 \qquad X^4/4!=\tfrac12E_{15}.
\]
E.3 rules out the asserted homomorphism.  All entries and identities
are verified by the program.

For the general bound, expand
\((I+N)^t=\sum_{k=0}^{r-1}\binom tkN^k\).
The condition on \(N^p\) implies
\(N^k\in p^{\lfloor k/p\rfloor}M_r(\mathbb Z_{(p)})\).
If \(k<p^2\), Legendre's formula gives
\(v_p(k!)=\lfloor k/p\rfloor\).  Since \(k!\binom tk\) has integral
coefficients, each summand is integral.  The group law holds over
\(\mathbb Q\) and hence over the torsion-free coefficient ring.

For sharpness, set \(d=p^2\) and give a rank-\(d+1\) shift the weights
\[
                  w_i=i/p^{\max(v_p(i)-1,0)}\quad(1\leq i\leq d).
\]
Every product of \(p\) consecutive weights is divisible by \(p\),
whereas \(\prod_{i=1}^{p^2}w_i=(p^2)!/p\).  Consequently
\(N^p\equiv0\pmod p\), but the leading coefficient in
\((I+N)^t\), namely \(N^{p^2}/(p^2)!\), has an entry \(1/p\).
For primes \(q\ne p\) the weights have the same \(q\)-valuations as
\(1,\ldots,p^2\), so all geometric fibres satisfy the required
one-parameter condition.  This also supplies genuine integral
examples at every prime.  The executed search through degree 12 gives
first failing degrees 4 and 9 for \(p=2,3\); no member fails for
\(p=5\) in that range, as the proven threshold 25 predicts.
\end{proof}

\begin{statement}{E.7}{Refuted}{Positive-characteristic and infinitesimal-base tests}
For every prime \(p\), the conventional unipotent group \(W_2/\mathbb F_p\)
contains \((1,0)\) of order \(p^2\), so it is not a value at 1 of a
\(\mathbb G_a\)-homomorphism.  More strongly, over
\(A=\mathbb Z/p^2\), the smooth affine group
\[
 G=\mathbb A^1_A,\qquad
 x*y=x+y+pC_p(x,y),\qquad
 C_p(x,y)=\frac{(x+y)^p-x^p-y^p}{p}
\]
and section \(g=1\) satisfy the manuscript's fibrewise hypothesis,
but admit no homomorphism \(\mathbb G_{a,A}\to G\) with value \(g\).
For \(p=2\) the law is \(x+y+2xy\), so the smallest ring in this
family is \(\mathbb Z/4\).  In every case \(g\) has order \(p\),
so this counterexample does not refute the positive-power assertion.
\end{statement}
\begin{proof}
For Witt vectors the addition law is
\[
 (a_0,a_1)+(b_0,b_1)=
 \left(a_0+b_0,\ a_1+b_1-
 \sum_{i=1}^{p-1}\frac1p\binom pi a_0^i b_0^{p-i}\right).
\]
The standard identity \([p](a_0,a_1)=(0,a_0^p)\) follows by adding
the ghost polynomials and reducing modulo \(p\); a second multiplication
by \(p\) gives zero.  Hence \((1,0)\) has order \(p^2\).  E.1 explains
exactly which proposed implication this refutes.  The program checks
the group law on all triples for \(p=2,3,5\), respectively 64, 729,
and 15,625 triples, and computes the three orders.

For the second example, \(C_p\) is an integral polynomial and an
additive two-cocycle, being the divided coboundary of \(x^p\).
Because \(p^2=0\) in \(A\), changing an argument of \(C_p\) by
a multiple of \(p\) does not change \(pC_p\).  The cocycle identity
therefore proves associativity of the displayed law.  The inverse
is \(-x-pC_p(x,-x)\).  Thus the underlying smooth affine line
does define a group.  Summing the cocycle along the \(p\) copies
of 1 gives \([p]g=p+p(p^p-p)/p=p^p=0\) in \(A\), while its
reduction has exact order \(p\).  Its geometric fibre is \(\mathbb G_a\),
and the section reduces to 1, so the hypothesis of scan 6 holds.
Suppose a polynomial \(f(t)\in A[t]\) defined the claimed homomorphism.
Then \(f(0)=0\), \(f(1)=1\), and
\[
       f(t+u)-f(t)-f(u)=pC_p(f(t),f(u)).          \tag{E.7.1}
\]
Modulo \(p\), \(f\) is an additive polynomial.  Let \(d\) be its
largest degree with coefficient not divisible by \(p\).  It exists
because \(f(1)=1\), and \(d\) is a power of \(p\).  The coefficient
of \(t^{(p-1)d}u^d\) on the right of (E.7.1) is
\(pa_d^p\ne0\) modulo \(p^2\).  On the left it is
\(\binom{pd}d a_{pd}\).  The coefficient \(a_{pd}\) is divisible
by \(p\) by the choice of \(d\), and \(\binom{pd}d\) is divisible
by \(p\) (indeed its valuation is 1).  The left coefficient
is therefore zero modulo \(p^2\), a contradiction.  This proves failure
in every degree.  Independently, the program tests all 1,024 candidate
polynomials over \(\mathbb Z/4\) of degree at most 6 with
\(f(0)=0\), \(f(1)=1\), using
the polynomial identity rather than only its values on the 16 pairs
of \(\mathbb Z/4\)-points; none succeeds.  It also checks the
deformed group laws on all triples of points for \(p=2,3,5\),
computes the orders of \(g\), and verifies the leading-term
obstruction for \(d=1,p,p^2,p^3\).
\end{proof}

\subsection{The doubtful lemma and the positive-power theorem}

\begin{statement}{E.8}{Proved}{Integer-valued polynomials and the finite-difference correction}
A polynomial \(P\in\mathbb Q[t]\) of degree at most \(d\) is
integer-valued on \(\mathbb Z\) if and only if
\[
             P(t)=\sum_{j=0}^d c_j\binom tj,\qquad c_j\in\mathbb Z.
\]
Integer values, even at every integer, do not imply integral monomial
coefficients.  For any commutative ring \(A\), a polynomial
\(P\in A[t_1,\ldots,t_r]\) which vanishes on every integral tuple
has all its coefficients annihilated by some positive integer.
If in addition \(P(0)=0\), some substitution
\(t_i\mapsto Dt_i\) makes it identically zero.
\end{statement}
\begin{proof}
Repeated forward differences give the Newton expansion
\(c_j=\Delta^jP(0)\).  Thus integral values imply integral coefficients
in the binomial basis; the converse follows because all binomial
polynomials take integral values at positive and negative integers.
This is the \(\mathbb Z\) case discussed in the supplied
\citet[\S1, p.~964]{E_Bhargava2009}; the preceding argument proves the
case we use directly, without invoking a general \(P\)-ordering theorem.
The polynomial \(\binom t2=(t^2-t)/2\) proves the warning.

For the assertion over \(A\), choose an individual degree bound \(d\)
in every variable.  Evaluation on \(\{0,\ldots,d\}^r\) is the tensor
product of integral Vandermonde matrices.  Multiplication by the
adjugate shows that a nonzero integer annihilates every coefficient
of a polynomial in its kernel.  Choose one such integer \(D\).
After replacing every variable by \(Dt_i\), each nonconstant
coefficient acquires a factor divisible by \(D\); the constant
coefficient was zero.  This proves the last assertion.

The degree-12 computation records every binomial basis element and
its monomial denominator.  It also checks the two specified matrix
families \(X=\log(1+J_{d+1})\) and the weighted family of E.6 for
each \(1\leq d\leq12\) and \(p=2,3,5\).  This is an exhaustive
search of those 72 cases, not an assertion of exhaustion of all
rational matrices.  The counterexamples and the proofs E.4--E.6
settle the implications for which the search was used.

Two further readings on C4 scan 13, fol.~654, can be tested exactly.
In (20), the reading \(X\in M_5(A)\) makes the displayed conclusion
\(X\in M_5(A)\) tautological.  With the alternative initial reading
\(X\in M_5(K)\), the conclusion is false: take
\(A=\mathbb Z_{(2)}\), \(K=\mathbb Q\), and
\(X=\tfrac12E_{12}\in M_5(K)\).  All its powers of degree at least
two vanish, so both signed higher-power combinations are integral,
although \(X\) is not.  This is a test of the literal alternatives,
not a proposed final transcription.  The cubic denominator \(3!\)
in the checked transcription is forced by the exponential addition
law; the first-pass \(6!\) alternative gives coefficient
\(-119/240\) of \(t^2uE_{14}\) in
\(P(t+u)-P(t)P(u)\), for a rank-four Jordan block.  Both calculations
are included in \texttt{E\_integrality.py}.  Finally, the phrase
``car.~0 i.e.\ a \(\mathbb Q\)-algebra'' on that page cannot be
an equivalence of ring-theoretic hypotheses: \(\mathbb Z\) is a
counterexample.  Our E.2 uses the stated \(\mathbb Q\)-algebra
hypothesis, and E.6 shows why it cannot simply be discarded.
\end{proof}

\begin{statement}{E.9}{Proved}{The affine dilation proposition of scans 14--16}
The proposition of \S50.2, fols.~654'--655', C4 scans 14--16,
has the following precise affine form.
Let \(A\) be a ring such that \(A\to A[1/n]\) is injective, and let
\(G=\operatorname{Spec}B\) be an affine group scheme of finite type
over \(A\).  Every homomorphism
\(\phi:\mathbb G_{a,A[1/n]}\to G_{A[1/n]}\) extends after precomposition
with multiplication by \(n^v\), for some \(v\geq0\), to a
homomorphism \(\psi:\mathbb G_{a,A}\to G\).
More generally, for an affine finite type \(A\)-scheme \(Y\), a
morphism \(\mathbb A^1_{A[1/n]}\to Y_{A[1/n]}\) admits such an
extension if and only if its value at zero extends to \(Y(A)\).
\end{statement}
\begin{proof}
Choose algebra generators \(b_1,\ldots,b_s\) of \(B\).  Their images
are polynomials \(\sum_j a_{ij}t^j\in A[1/n][t]\).  Their constant
coefficients are in \(A\), because a group homomorphism takes zero
to the identity.  One exponent \(v\) makes
\(n^{vj}a_{ij}\in A\) for all \(j\geq1\) and all \(i\).  The
substituted polynomials define the required algebra map: every
relation vanishes after localization and therefore vanishes in
\(A[t]\), by the injectivity assumption.  The two group-law maps
agree after localization, so the same argument in \(A[t,u]\) proves
that the extension is a group homomorphism.  For \(Y\), the same
argument applies to its coordinate ring once its constant terms
come from the specified zero section; evaluating at zero proves
necessity.  Finite affine covers give the corresponding statement
for a quasi-compact base with schematically dense \(D(n)\), since
the extensions are unique and hence glue.

This proves the proposition actually written in C4 scans 14--16.
The substitution is \(t\mapsto n^vt\), not \(t\mapsto t^{n^v}\).
The latter is not additive in characteristic zero.  The correction
is supported both by the type of the map and by formula (22) in
scan 14.
\end{proof}

\begin{statement}{E.10}{Proved}{A complete positive-power theorem for linear group schemes}
Let \(A\) be any commutative ring and
\(G\hookrightarrow\operatorname{GL}_{r,A}\) a closed immersion of
affine group schemes of finite presentation over \(A\).  A section
\(g\in G(A)\) is fibrewise quasi-unipotent in the manuscript's sense
if and only if there exist an integer \(m>0\) and a homomorphism
\(\phi:\mathbb G_{a,A}\to G\) with \(\phi(1)=g^m\).
This proves the ``Th\'eor\`eme (!)'' of fol.~652', C4 scans 10--11,
under the specified linearity hypothesis, including all linear
examples over \(\mathbb Z\), \(\mathbb Z_p\), and \(\mathbb Z/p^2\).
\end{statement}
\begin{proof}
Only necessity needs proof.  For each prime ideal of \(A\), a
positive power of the matrix \(g\) is unipotent in the residue
field.  The sets
\[
 V\bigl(\hbox{entries of }(g^a-I)^r\bigr),\qquad a=1,2,\ldots,
\]
cover \(\operatorname{Spec}A\).  Each is open and closed in the
constructible topology, which is compact; alternatively, apply
compactness to the Boolean algebra of finitely generated radical
conditions.  For clarity, this compactness holds for arbitrary
\(A\): an ultrafilter \(\mathcal U\) on \(\operatorname{Spec}A\)
defines the prime ideal
\(\mathfrak p=\{f:V(f)\in\mathcal U\}\).
The identities for \(V(fg)\), and the inclusion
\(V(f)\cap V(g)\subset V(f+g)\), prove primality and additivity;
the ultrafilter converges to \(\mathfrak p\) for the topology
generated by the clopen sets \(V(f),D(f)\).  Thus this topology
is compact.  A finite subcover therefore exists.  Taking the least
common multiple of its exponents shows that \(N=g^a-I\) is
nilpotent on every residue field.  As in the Cayley--Hamilton
argument of E.2, the finitely many characteristic coefficients
generate a nilpotent ideal, and \(N^b=0\) for some \(b\).

Choose \(d=(b-1)!\), with \(d=1\) if \(b=1\).  Every polynomial
\(\binom{dt}j\), \(1\leq j<b\), has integral coefficients: its
constant term is zero, and each other coefficient contains a
factor \(d\), which is divisible by \(j!\).  Thus
\[
                   P(t)=\sum_{j=0}^{b-1}\binom{dt}jN^j
\]
belongs to \(M_r(A[t])\).  The binomial identity proves
\(P(t+u)=P(t)P(u)\) over the universal torsion-free ring
\(\mathbb Z[z]/(z^b)\), and hence after every base change.
Its inverse is \(P(-t)\), and \(P(n)=g^{adn}\) for every integer
\(n\).  It therefore defines a homomorphism to \(\operatorname{GL}_r\).

There are finitely many defining equations for \(G\) inside
\(\operatorname{GL}_r\).  Their pullbacks by \(P(t)\), including
the inverse-determinant coordinate, vanish at every integer and
at zero.  E.8 supplies one further dilation \(t\mapsto Dt\)
which makes every pullback identically zero.  The resulting
homomorphism factors through \(G\), and its value at 1 is
\(g^{adD}\).  Sufficiency is the definition of quasi-unipotence.
\end{proof}

\begin{statement}{E.11}{Open}{The remaining non-linear scope of scans 6 and 10}
We have not proved the following full-source statement:
for every quasi-compact scheme \(S\), every separated group scheme
\(G/S\) locally of finite presentation, and every fibrewise
quasi-unipotent \(g\in G(S)\), a positive power of \(g\) is the
value at 1 of a homomorphism \(\mathbb G_{a,S}\to G\).
Likewise E.2 has not established the assertion for all group
schemes over all \(\mathbb Q\)-algebras without a linearity hypothesis.
\end{statement}
\noindent
Our proved cases are E.2, E.3, E.9, and E.10.  The executed matrix
and polynomial tests are E.4--E.8, with their complete outputs.
The precise missing step is extension from a general affine or
non-affine group scheme to a faithful finite locally free linear
representation, or a substitute for that step.  We do not assume
such a representation exists.  A first remaining quantified test
in base dimension zero is \(S=\operatorname{Spec}
\mathbb Q[\epsilon]/(\epsilon^2)\), with arbitrary separated affine
finite presentation \(G/S\) and no assumed closed linear embedding;
the analogous integral test is over \(\mathbb Z/4\).
These are cases not covered by the proof, not claims that the
corresponding questions are open in the literature.  The example
E.7 does not obstruct the positive-power theorem, because its
section has order \(p\).  This distinction is essential to the status
assigned to the full manuscript assertion.

\subsection{The envelope as it is actually defined}

\begin{statement}{E.12}{Proved}{The terminology and the finite-stage pushout}
The ``m\'etab\'elien'' groups of \S50.4, fol.~660, C4 scan 25,
are finitely generated virtually abelian groups.  Let \(\Gamma\)
be such a group and \(M\simeq\mathbb Z^r\) a normal subgroup of
finite index, with \(H=\Gamma/M\).  Choose a set section
\(s:H\to\Gamma\), \(s(1)=1\), and write its action and cocycle as
\[
 a_h(m)=s(h)ms(h)^{-1},\qquad
 c(h,k)=s(h)s(k)s(hk)^{-1}\in M.
\]
There is a smooth affine \(\mathbb Z\)-group scheme
\[
 \underline\Gamma_M=\coprod_{h\in H}\mathbb A^r_{\mathbb Z},\qquad
 \mathcal O(\underline\Gamma_M)=\prod_{h\in H}\mathbb Z[x_1,\ldots,x_r],
\]
whose multiplication on the \((h,k)\) block is
\[
               (m,h)(n,k)=(m+a_hn+c(h,k),hk).       \tag{E.12.1}
\]
It lies in an exact sequence of fppf group sheaves
\(1\to\mathbb G_a^r\to\underline\Gamma_M\to H_{\mathbb Z}\to1\),
and \(\Gamma\to\underline\Gamma_M(\mathbb Z)\) is an isomorphism.
The construction is independent, up to the unique coordinate change
specified by the two sections, of the chosen section.
\end{statement}
\begin{proof}
The identity
\(c(h,k)+c(hk,l)=a_hc(k,l)+c(h,kl)\) follows by associating
\(s(h)s(k)s(l)\) in two ways.  It proves associativity of (E.12.1);
the identity is \((0,1)\), and solving the same formula gives
polynomial inverses.  Each block is an affine space over
\(\mathbb Z\), proving the stated geometric properties.  Projection
to \(H\) has a section on every component as a morphism of schemes,
so the sequence is exact even on sections over an arbitrary base
scheme, with locally constant \(H\)-labels.  Replacing \(s(h)\) by
\(b(h)s(h)\), \(b(h)\in M\), replaces the vector coordinate by
\(m-b(h)\); this explicitly identifies the two laws.

For a disconnected test scheme \(T\), the label is a locally
constant function \(T\to H\), not necessarily one fixed element
of \(H\).  This is exactly the correction made by Zariski
sheafification in C4 scans 18--20.  Formula (E.12.1) is also the
pushout of the extension \(1\to M\to\Gamma\to H\to1\) along
\(M_{\mathbb Z}\to\mathbb G_a^r\), with the prescribed conjugation
action.  This proves its universal property among maps of such
extensions directly from the cocycle formula.
\end{proof}

\begin{statement}{E.13}{Proved}{Coordinate ring and fibres of the full envelope}
Set \(M_n=nM\), \(H_n=\Gamma/nM\), and
\(E_\Gamma=\varprojlim_n\underline\Gamma_{nM}\), where indices are
ordered by divisibility.  Then
\[
 \mathcal O(E_\Gamma)=
 \varinjlim_n\prod_{h\in H_n}\mathbb Z[x_1,\ldots,x_r].             \tag{E.13.1}
\]
For \(n\mid n'\), the coordinate changes on each component are
\(x\mapsto (n'/n)x+b\), with explicitly determined
\(b\in\mathbb Z^r\).  The construction is independent of the
chosen normal lattice.  Its kernel over the proconstant group
\(\widehat\Gamma_{\mathbb Z}=\varprojlim H_{n,\mathbb Z}\) is
\[
 E^0_\Gamma\simeq
 \operatorname{Spec}\bigl(\mathbb Z+(x_1,\ldots,x_r)
                              \mathbb Q[x_1,\ldots,x_r]\bigr).    \tag{E.13.2}
\]
For every prime \(p\), there is a canonical isomorphism
\[
                   (E_\Gamma)_{\mathbb F_p}\simeq
                   \widehat\Gamma_{\mathbb F_p}.                 \tag{E.13.3}
\]
In particular these fibres are reduced, of Krull dimension zero,
and generally not of finite type.  The finite-stage scheme of
E.12 instead has \(|H|\) smooth components of dimension \(r\) in
every characteristic.
\end{statement}
\begin{proof}
The inclusions \(n'M\hookrightarrow nM\) and identity of \(\Gamma\)
give the transition maps by E.12.  They are affine, so the inverse
limit is represented by the indicated direct limit of rings.
The subgroups \(nM\) are cofinal among all finite-index subgroups
of \(M\), and among normal sublattices used in the construction.
Any two normal lattices have a common finite-index intersection;
this proves independence.  Every finite-index normal subgroup of
\(\Gamma\) contains some \(nM\), so \(\varprojlim H_n\) really is
the full profinite completion.

On the identity components the transition is multiplication by
\(n'/n\).  Its coordinate-ring limit is the union of
\(\mathbb Z[x_1/n,\ldots,x_r/n]\) in \(\mathbb Q[x_1,\ldots,x_r]\).
This union is precisely (E.13.2): clear the finitely many rational
denominators of every nonconstant monomial; the constant term
must be integral.  This proves the kernel assertion as a statement
about affine group schemes.

In rank one its points are coherent sequences
\((a_n)\), \(a_n=m a_{mn}\), as in scan 29, fol.~662,
formula (56).  If an integer \(q>0\) annihilates \(A\),
every \(a_n=q a_{qn}=0\).  If every residue characteristic
is positive, such a \(q\) does exist: otherwise
\(A\otimes\mathbb Q\ne0\) has a prime ideal, producing a
characteristic-zero residue field of \(A\).
For a \(\mathbb Q\)-algebra the sequence is uniquely
\(a_n=a_1/n\), so the point group is \((A,+)\).
Characteristic zero alone is insufficient: for
\(A=\mathbb Z\) every coherent sequence is zero.
Also, ``all residue characteristics divide \(q\)'' does not
imply \(qA=0\) for that same \(q\), as
\(A=\mathbb Z/4\), \(q=2\), shows.  These give the precise
corrections to the bracketed remarks in scan 29.

After reduction modulo \(p\), the transition from the \(np\)-stage
to the \(n\)-stage has zero linear part.  Hence every polynomial
at stage \(n\) becomes a locally constant function at stage \(np\).
Conversely all locally constant functions on each finite \(H_n\)
are already in the stage coordinate ring.  Therefore its direct
limit is exactly
\(\varinjlim_n\mathbb F_p^{H_n}\), the coordinate ring of the
proconstant group.  It is a filtered union of products of fields,
so is reduced and absolutely flat; every prime is maximal.

For \(\Gamma=\mathbb Z\), (E.13.1) can also be described without
chosen coordinates: it is the ring of integer-valued rational
quasipolynomials on \(\mathbb Z\), namely functions which are
rational polynomials on the classes of one finite modulus.
An integer-valued polynomial on a class has integral monomial
coefficients after refining the modulus enough to clear all
nonconstant denominators; its constant term is its integral value
at the chosen representative.  This proves the description in
both directions.  It contains nontrivial periodic idempotents,
which the divided-power polynomial algebra alone does not contain.
\end{proof}

\begin{statement}{E.14}{Proved}{The universal property for linear targets, with the correct colimit}
For any ring \(A\), any finitely generated virtually abelian
\(\Gamma\), and any closed finite presentation subgroup scheme
\(G\hookrightarrow\operatorname{GL}_{r,A}\), restriction along
\(\Gamma_A\to(E_\Gamma)_A\) is a bijection
\[
 \operatorname{Hom}_{A\text{-grp}}((E_\Gamma)_A,G)
 \simeq
 \{f:\Gamma\to G(A)\mid f\text{ is fibrewise quasi-unipotent}\}.   \tag{E.14.1}
\]
The left side is
\(\varinjlim_n\operatorname{Hom}_{A\text{-grp}}
((\underline\Gamma_{nM})_A,G)\), not the inverse limit printed in
formula (58), C4 scan 31.  This proves the universal assertion of
\S50.5, fol.~663, scan 32, for the stated targets.
\end{statement}
\begin{proof}
An affine morphism from a filtered inverse limit to a finitely
presented affine scheme factors through a finite stage: choose
images for finitely many algebra generators and then a later
stage where their finitely many relations vanish.  Applied also
to the multiplication and inverse equations, this gives the
displayed colimit of group homomorphisms.  This elementary algebra
argument explains the direction of the limit correction.

Let \(f\) be on the right of (E.14.1).  Choose a basis of the normal
lattice \(M\).  Applying E.10 to its finitely many images and
taking a common multiple of the resulting exponents gives
one-parameter homomorphisms with values \(f(ne_i)\) at 1.
We take the explicit matrix polynomials of E.10, which are
polynomials in powers of \(f(e_i)\).  Since \(M\) is abelian,
these homomorphisms commute with each other and with every
\(f(m)\), \(m\in M\), over \(A\) itself.  Their product is a vector-group homomorphism
\(v:(nM)\otimes\mathbb G_a\to G\) extending \(f|_{nM}\).

For each of the finitely many coset representatives of \(\Gamma/M\),
the required identity between conjugation of \(v\) and the
prescribed linear action on its source is again a finite set of
polynomial identities.  They hold on every integral tuple.
Another application of E.8, and a further common dilation,
make them identities over \(A\).  The cocycle formula E.12 now
extends \(f\) to \(\underline\Gamma_{nM}\).

If two extensions agree on \(\Gamma\), pass to a common stage.
On each of its finitely many affine-space components their
coordinate differences vanish on the integral lattice.  E.8
kills every such difference after a common refinement, proving
uniqueness in the colimit.  Finally, any map from a stage has
quasi-unipotent restriction: a positive power of each element of
\(\Gamma\) lies in the lattice, and its image lies on the image of
the line through the corresponding integral lattice vector.
This proves the reverse inclusion and the bijection.
\end{proof}

\noindent
For the full reading ``\(G\) affine over \(S\)'' or the stronger
reading ``affine over \(S_0\)'' in scan 32, E.11 records the
remaining linearization issue.  The computations E.4--E.8 are the
attack on the section-extension step, and E.14 explains exactly
which proved correction is sufficient for the universal property.
We do not turn that unproved generality into a claim about all
affine group schemes.

\subsection{The explicit modular example and its points}

\begin{statement}{E.15}{Proved}{The lattice and quotient in the modular example}
Fix
\[
 \rho=\begin{pmatrix}0&1\\-1&1\end{pmatrix},\qquad
 \sigma=\begin{pmatrix}0&-1\\1&0\end{pmatrix},\qquad
 \lambda_0=-\sigma\rho\sigma\rho,
 \quad\lambda_1=-\rho\sigma\rho\sigma.
\]
Then \(\Pi_0=\langle\lambda_0,\lambda_1\rangle\) is a normal free
subgroup of rank two and index 12 in \(\operatorname{SL}_2(\mathbb Z)\).
For
\(\mathfrak S=\operatorname{SL}_2(\mathbb Z)/[\Pi_0,\Pi_0]\),
\[
 1\longrightarrow M=\Pi_0^{\mathrm{ab}}\simeq\mathbb Z^2
 \longrightarrow\mathfrak S\longrightarrow H\longrightarrow1,
 \quad
 H=\langle r,s\mid r^6=1,\ s^2=r^3,\ srs^{-1}=r^{-1}\rangle
\]
has \(|H|=12\).  The action on the basis \((\lambda_0,\lambda_1)\)
is
\[
             R=\begin{pmatrix}0&-1\\1&-1\end{pmatrix},\qquad
             S=\begin{pmatrix}0&1\\1&0\end{pmatrix}.
\]
Moreover \(\mathfrak S^{\mathrm{ab}}\simeq\mathbb Z/12\).
\end{statement}
\begin{proof}
Matrix multiplication gives
\[
 \rho\lambda_0\rho^{-1}=\lambda_1,\qquad
 \rho\lambda_1\rho^{-1}=(\lambda_1\lambda_0)^{-1},\qquad
 \sigma\lambda_0\sigma^{-1}=\lambda_1,\qquad
 \sigma\lambda_1\sigma^{-1}=\lambda_0.
\]
Thus the subgroup generated by the two \(\lambda_i\) is normal.
Using the amalgam presentation
\(\operatorname{SL}_2(\mathbb Z)=C_4*_{C_2}C_6\), imposing
\(\lambda_0=1\) gives \(\sigma\rho\sigma^{-1}=\rho^{-1}\);
\(\lambda_1=1\) is its conjugate.  The resulting presentation is
the one displayed for \(H\).  Every word reduces to
\(r^is^j\), \(0\leq i<6\), \(0\leq j<2\), so its order is at
most 12.  The calculation in \texttt{E\_envelope.py} realizes
12 distinct classes as
\[
 H\simeq\operatorname{SL}_2(\mathbb Z/4)/
                 \langle\overline\lambda_0,\overline\lambda_1\rangle;
\]
the denominator is a normal Klein four group.  This proves the
order and the index without a coset-order guess.

The subgroup \(\Pi_0\) intersects each conjugate of the amalgam
factors trivially, because the factors of orders 4 and 6 inject
into \(H\).  It therefore acts freely on the amalgam tree.  Its
quotient graph has \(12/4+12/6=5\) vertices and \(12/2=6\) edges,
and is connected.  Its fundamental group is free of rank
\(6-5+1=2\).  The two generating elements therefore give a basis
of its abelianization.  Abelianizing the four conjugation formulas
gives \(R,S\).  Finally, quotienting by \([\Pi_0,\Pi_0]\), which
is contained in the commutator subgroup of \(\operatorname{SL}_2(\mathbb Z)\),
does not change its abelianization.  Abelianizing the amalgam
presentation gives
\(\langle r,s\mid6r=0,4s=0,3r=2s\rangle\simeq\mathbb Z/12\).
These are the discrete groups in \S51.1, fol.~664, C4 scan 33;
the word ``profinis'' in that passage is incompatible with the
uncapped lattices in the same displayed sequence.  If the deleted
hats are retained instead, the sequence must be read as one of
profinite completions; it is then a different typed assertion.
Our input is the discrete group specified in the question.
\end{proof}

\begin{statement}{E.16}{Proved}{An explicit integral coordinate ring for the modular finite stage}
The group scheme \(\underline{\mathfrak S}_M\) of \S51.1 has
\[
       \mathcal O(\underline{\mathfrak S}_M)
           =\prod_{h\in H}\mathbb Z[x,y].                       \tag{E.16.1}
\]
Its multiplication is (E.12.1) with the following completely
specified cocycle.  Take the twelve representatives
\[
 1,r,s,r^2,rs,sr,s^2,r^2s,rs^2,sr^2,s^3,r^2s^2.
\]
Associate to a word its linear action obtained from \(R,S\), and
its translation obtained by composing
\[
             r\mapsto(R,0),\qquad
             s\mapsto\bigl(S,(1/3,-1/3)^t\bigr).
\]
If \(t_h\) is that translation, set
\[
                    c(h,k)=t_h+a_ht_k-t_{hk}.                    \tag{E.16.2}
\]
Every \(c(h,k)\) is integral.  The model is smooth and flat of
relative dimension two over \(\mathbb Z\), with twelve reduced
geometrically connected components in each fibre, including 2
and 3.  Its number of \(\mathbb F_p\)-points is \(12p^2\).
\end{statement}
\begin{proof}
The affine transformations assigned to \(\rho,\sigma\) satisfy
\(\rho^6=\sigma^4=1\), \(\rho^3=\sigma^2\), and send
\(\lambda_0,\lambda_1\) to translations by \((1,0)\) and
\((0,1)\).  Hence the translation of a word in \(\Pi_0\) is
precisely its abelianized lattice coordinate.  The difference
word defining \(c(h,k)\) lies in \(\Pi_0\); this proves integrality
and shows that this is the cocycle of the specified discrete
extension.  Thus E.12 proves the coordinate ring and geometric
claims.  The complete multiplication table, all 144 cocycle
values, and all 1,728 cocycle identities are independently
computed in \texttt{E\_envelope.json}; the program also verifies
the four integral matrix conjugation identities from E.15.

For \(p=2,3,5,7,11,13\), the point counts are respectively
\(48,108,300,588,1452,2028\).  The discrete quotients
\(\mathfrak S/nM\), generated by the images of \(\rho,\sigma\),
have computed orders \(12,48,108,192,768,972\) for
\(n=1,2,3,4,8,9\), agreeing with \(12n^2\).
The denominators 3 in the auxiliary translations do not introduce
singularities: the group law itself uses only the integral
cocycle (E.16.2).
\end{proof}

\begin{statement}{E.17}{Proved}{Integral, local, and profinite-ring points of the envelope}
For every finitely generated virtually abelian \(\Gamma\),
\[
 E_\Gamma(\mathbb Z)=\Gamma,\qquad
 E_\Gamma(\mathbb Z_p)=\widehat\Gamma,
 \qquad
 E_\Gamma(\widehat{\mathbb Z})=\prod_p\widehat\Gamma.              \tag{E.17.1}
\]
The local equality is an isomorphism of topological groups, with
the inverse-limit topology on the left.  For the finite-stage
model one instead has
\[
 \underline\Gamma_M(\widehat{\mathbb Z})
       =\prod_p\underline\Gamma_M(\mathbb Z_p),\qquad
 \underline\Gamma_M^{\Delta}(\widehat{\mathbb Z})
       \simeq\widehat\Gamma,                                  \tag{E.17.2}
\]
where \(\Delta\) means that all local images in \(H\) coincide.
These statements apply in particular to \(\mathfrak S\) of E.15.
\end{statement}
\begin{proof}
The map to the first stage identifies a \(\mathbb Z\)-point with
an element of \(\Gamma\).  At stage \(n\), compatibility says
that its lattice coordinate equals a chosen residue modulo
\(nM\) plus \(n\) times an integral coordinate.  This determines
both the residue and the coordinate, so no extra point lies
over a given element of \(\Gamma\).  The natural map from
\(\Gamma\) supplies all the points.

The kernel (E.13.2) has no nonzero \(\mathbb Z_p\)-points:
its coordinates form systems \(x_n=mx_{mn}\), so each fixed
\(x_n\) is divisible by every power of \(p\) and is zero.
Each finite-stage group \(\underline\Gamma_{nM}(\mathbb Z_p)\)
is profinite, being a finite disjoint union of copies of
\(\mathbb Z_p^r\).  The map from \(\Gamma\) to it extends uniquely
and continuously to \(\widehat\Gamma\).  These extensions are
compatible in \(n\) and give a map
\(\widehat\Gamma\to E_\Gamma(\mathbb Z_p)\) inverse to projection
on the proconstant quotient: their composite is the identity
on the dense subgroup \(\Gamma\), and the other composite
differs from the identity by the zero kernel just calculated.

For an affine scheme \(X=\operatorname{Spec}B\),
\(\operatorname{Hom}(B,\prod_p\mathbb Z_p)
 =\prod_p\operatorname{Hom}(B,\mathbb Z_p)\).
This gives the third equality of (E.17.1) and the first in
(E.17.2), without a connectedness assumption on
\(\operatorname{Spec}\widehat{\mathbb Z}\).  In the finite-stage
model, requiring a common component \(h\in H\) leaves exactly
one copy of \(M\otimes\widehat{\mathbb Z}\) above every \(h\),
with the same cocycle as the profinite completion of \(\Gamma\).
The topology on \(M\) induced from \(\Gamma\) is its full
profinite topology because all \(nM\) are normal finite-index
subgroups.  This proves (E.17.2).

Thus the \(\Delta\) in C4 scan 35, formula (13), is necessary.
For the full envelope, requiring a common value in every finite
quotient across all primes similarly cuts out the diagonal
\(\widehat\Gamma\subset\prod_p\widehat\Gamma\).
An assertion of equality with all
\(\widehat{\mathbb Z}\)-points would be false.
\end{proof}

\subsection{Comparisons with the proposed hulls}

\begin{statement}{E.18}{Refuted}{Lead 3 does not identify the envelope with a pro-algebraic or Malcev hull}
Let \(P_{\mathbb Q}(\Gamma)\) denote the pro-algebraic group
universal for all finite-dimensional rational representations
of a discrete group \(\Gamma\), and let \(U_{\mathbb Q}(\Gamma)\)
denote its unipotent completion.  For finitely generated
virtually abelian \(\Gamma\), there are canonical quotient maps
\[
 P_{\mathbb Q}(\Gamma)\longrightarrow(E_\Gamma)_{\mathbb Q}
                 \longrightarrow U_{\mathbb Q}(\Gamma),         \tag{E.18.1}
\]
defined by the natural maps from \(\Gamma\).  Both can be strict.
For \(\Gamma=\mathbb Z\),
\[
 (E_{\mathbb Z})_{\mathbb Q}\simeq
             \mathbb G_{a,\mathbb Q}\times\widehat{\mathbb Z}_{\mathbb Q},
 \qquad U_{\mathbb Q}(\mathbb Z)=\mathbb G_{a,\mathbb Q}.
\]
The character \(n\mapsto2^n\) belongs to the first theory and
does not factor through \((E_{\mathbb Z})_{\mathbb Q}\).
For \(\Gamma=\mathfrak S\), the unipotent completion is trivial,
but its envelope retains the quotient \(\mathfrak S_3\).
\end{statement}
\begin{proof}
Every representation of the envelope restricts to a
quasi-unipotent representation by E.14.  Conversely every such
representation extends by the same theorem.  This identifies
its representation category as the full tensor subcategory of
quasi-unipotent representations of \(\Gamma\), giving the first
map in (E.18.1) by the defining universal property of
\(P_{\mathbb Q}\).  Its coordinate ring is generated by matrix
coefficients of these representations, so the map is a
pro-algebraic quotient.  Unipotent representations are
quasi-unipotent, giving the second map in exactly the same way.

At the \(n\)-stage for \(\mathbb Z\), a point in residue class
\(a\) has additive coordinate \(a+nt\).  Over \(\mathbb Q\)
this gives an isomorphism with
\(\mathbb G_a\times(\mathbb Z/n)\), and the transition preserves
the first coordinate and reduces the second.  Taking limits
proves the displayed decomposition.  Its finite-dimensional
representations have only roots of unity as semisimple
eigenvalues of the distinguished generator; \(2\) is not one.

For \(\mathfrak S\), its finite abelianization gives
\(\operatorname{Hom}(\mathfrak S,\mathbb G_a(\mathbb Q))=0\).
If a homomorphism to a unipotent group had nontrivial Zariski
closure, that nontrivial unipotent group in characteristic zero
would have a nontrivial vector-group abelianization.  Composing
with a nonzero linear functional would contradict the displayed
vanishing.  This can also be seen directly by taking the first
nonzero superdiagonal in an upper-unitriangular realization.
Thus the unipotent completion is trivial.  Reduction modulo 2
of the quotient \(H\) supplies the nontrivial finite quotient
\(\mathfrak S_3\) of the envelope.

The integral Hall model of the rank-one nilpotent lattice is
\(\mathbb G_{a,\mathbb Z}\); the Lazard model of its additive
pro-\(p\) completion is the additive \(\mathbb Z_p\)-Lie lattice.
On this abelian example there are no Baker--Campbell--Hausdorff
denominators.  The maps to these models are the first-stage
projection \(E_{\mathbb Z}\to\mathbb G_a\) and its base change
or map on \(\mathbb Z_p\)-points.  They are not isomorphisms:
the envelope has every finite constant cyclic quotient, whereas
\(\mathbb G_a\) has no nontrivial morphism to a finite constant
group over \(\mathbb Q\).  This is a separating example for the
proposed identification with an integral Malcev or Lazard--Hall
form, without pretending that those constructions are a single
canonical integral functor on all virtually abelian groups.

There is also a direct coordinate-ring distinction from divided
powers.  In rank one the usual divided-power polynomial algebra is
\(D=\bigoplus_{j\geq0}\mathbb Z\,t^j/j!\subset\mathbb Q[t]\).
Its multiplication and coproduct are
\[
 t^{[i]}t^{[j]}=\binom{i+j}{i}t^{[i+j]},\qquad
 \Delta(t^{[n]})=\sum_{i+j=n}t^{[i]}\otimes t^{[j]}.
\]
It is not \(\mathcal O(E^0_{\mathbb Z})=\mathbb Z+t\mathbb Q[t]\):
the latter contains \(t/2\), whereas the degree-one part of
\(D\) is \(\mathbb Zt\).  Inclusion gives a specified morphism
\(E^0_{\mathbb Z}\to\operatorname{Spec}D\).
Their characteristic-\(p\) fibres also differ: the first is a
point, while the class of \(t\) is nonzero in \(D/pD\) and has
\(t^p=0\).  Thus neither the coordinate ring of the vector group
nor its familiar divided-power distribution algebra can silently
be identified with the envelope's kernel.  E.3 gives the exact
role divided powers do play in the section-extension problem.

For Hain's completion relative to a finite quotient
\(\rho:\Gamma\to H\), with \(\Gamma\) virtually abelian, there
is a canonical compatible map
\((E_\Gamma)_{\mathbb Q}\to\mathcal R_\rho(\Gamma)\).
Indeed every representation of the latter has the image of
some power of each \(\gamma\) in a unipotent group; E.14
extends every finite-dimensional quotient representation, and
uniqueness makes the extensions compatible.  For general
reductive quotients such a comparison can fail in both
directions.  Take \(\Gamma=\mathbb Z\) and
\(\rho(n)=2^n\in\mathbb G_m(\mathbb Q)\).  Its relative
completion is \(\mathbb G_a\times\mathbb G_m\), with
\(n\mapsto(n,2^n)\).  To prove this, the closure of a cyclic
image in any candidate extension is commutative.  Multiplicative
Jordan decomposition splits it into a diagonalizable part and
a unipotent part.  The former maps isomorphically onto
\(\mathbb G_m\), since its kernel must also be unipotent;
the latter is generated by one logarithm, so is the image of
\(\mathbb G_a\).  This gives the required unique map.
Uniqueness can also be checked by density of
\(\{(n,2^n):n\in\mathbb Z\}\): in a nonzero Laurent polynomial
\(\sum_jP_j(n)2^{jn}\), the largest \(j\) with nonzero
\(P_j\) dominates for positive integers \(n\), so the expression
cannot vanish identically.

There is no compatible map from \((E_{\mathbb Z})_{\mathbb Q}\)
to this relative completion, because its multiplicative
coordinate would extend the excluded character \(2^n\).
There is no compatible map in the other direction either:
the connected group \(\mathbb G_a\times\mathbb G_m\) has no
nontrivial morphism to a constant finite group, whereas the
envelope's quotient \(\mathbb Z/2\) is nontrivial on the
distinguished discrete subgroup.  Both receive the canonical
quotient map from the full pro-algebraic completion.
\end{proof}

\begin{statement}{E.19}{Proved}{The relative completion with abelian kernel is the finite stage}
Let \(\Gamma_0=\operatorname{SL}_2(\mathbb Z)\), with the specified
quotient \(\Gamma_0\to H\) of order 12 and kernel \(\Pi_0\).
In characteristic zero, the universal pro-algebraic extension
of \(H\) by a commutative unipotent group receiving this map is
\((\underline{\mathfrak S}_M)_{\mathbb Q}\).  Equivalently, if
\(\mathcal R\) is relative unipotent completion for
\(\Gamma_0\to H\), and \(\mathcal U\) its unipotent kernel, then
\[
       \mathcal R/\overline{[\mathcal U,\mathcal U]}
              \simeq (\underline{\mathfrak S}_M)_{\mathbb Q}.     \tag{E.19.1}
\]
This group is not \((E_{\mathfrak S})_{\mathbb Q}\): its component
group is the finite group \(H\), while the latter has the
infinite proconstant quotient \(\widehat{\mathfrak S}\).
The same statement holds for \(\operatorname{GL}_2(\mathbb Z)\)
and its quotient of order 24, with the orientation-reversing
action included.
\end{statement}
\begin{proof}
We use the universal property of relative completion formulated
in \citet[\S1, pp.~47--48]{E_Hain1998}; the abelian-kernel
calculation can be proved directly.  A map to an extension of
\(H\) by a commutative unipotent group sends \(\Pi_0\) to that
commutative kernel, so it kills \([\Pi_0,\Pi_0]\).  Over
\(\mathbb Q\), a commutative unipotent group is a vector group:
finite logarithm and exponential identify it with its Lie
algebra.  The induced homomorphism from
\(M=\Pi_0^{\mathrm{ab}}\) therefore extends uniquely and
linearly to \(M\otimes\mathbb G_a\).  The conjugation and
cocycle identities extend by linearity.  E.12 consequently
extends the original map uniquely to
\((\underline{\mathfrak S}_M)_{\mathbb Q}\).

Conversely that finite-stage group is itself an extension of
\(H\) by a vector group and receives the prescribed map.
Thus it has exactly the universal property of the quotient on
the left of (E.19.1).  This argument also works for inverse
limits of vector groups, so it gives the pro-algebraic assertion.
The difference of component groups proves non-isomorphism with
the full envelope.  The proof for the order-24 quotient is
identical: its normal lattice is still \(\Pi_0^{\mathrm{ab}}\),
but the additional generator
\(\tau'=\left(\begin{smallmatrix}0&1\\1&0\end{smallmatrix}\right)\)
conjugates \(\rho,\sigma\) to their inverses and sends
\(\lambda_0\mapsto\lambda_1^{-1}\),
\(\lambda_1\mapsto\lambda_0^{-1}\).  Its action on \(M\) is
\(-S\).  The program verifies these identities and the order
\(96/4=24\) of the corresponding quotient of
\(\operatorname{GL}_2(\mathbb Z/4)\).
\end{proof}

\begin{statement}{E.20}{Open}{The remaining comparison with the named automorphism schemes}
E.15--E.19 compute the group \(\mathfrak S\), its finite-stage
model, its full envelope, their fibres, their points, and the
relative completion with abelian kernel.  They do not identify
either of these group schemes with all versions of
\(M_{\rho,\sigma}\) or \(N_{\rho,\sigma}\) in \S49--\S51.
E.24 below computes the fine parameter scheme of C4 scan 41
over \(\mathbb Z\), and identifies the actual automorphism functor
with both conjugator signs positive over \(\mathbb Q\).
E.25 disproves the unqualified assertion in scan 36 that the two
generator images determine an automorphism over every base.
The exact remaining task is to compare every signed version of
\(M_{\rho,\sigma}\) and \(N_{\rho,\sigma}\) with the functor of
automorphisms of \(\underline{\mathfrak S}_M\) satisfying
\S51.2, conditions (15)--(23), specifying whether linearity on
\(\underline M\) is imposed, whether the conjugators and their
signs are remembered, and which fppf quotient removes redundant
parameters.
The group-scheme half of the After BISITE Fermat-tower follow-up
also asks for comparisons relative to
\((\mathbb Z/n)^2\) and
\(\mathfrak S_3\times\widehat{\mathbb Z}^{\,*}\).
The order-12/order-24 calculation in E.19 does not establish
these comparisons.
\end{statement}
\noindent
The first still unproved identification with the named signed
functors in the requested residue-characteristic list is
\(A=\mathbb F_2\), with both locally existing conjugators and
the sign data included.  We do not claim this case is open in
the literature.  The input data are fully explicit:
the integral cocycle (E.16.2), the matrices \(R,S\), and the
rank-two vector group.  Besides the finite quotient and cocycle
checks in \texttt{E\_envelope.py}, the program
\texttt{E\_automorphisms.py} computes the integral parameter
scheme, its points in characteristics 2 and 3 and over
\(\mathbb Z/4,\mathbb Z/8,\mathbb Z/9\), all sign alternatives
in scan 41, the generic automorphism composition, and the
infinitesimal obstruction.  The source itself introduces
\(\underline Z_{\rho,\sigma}\) as a subgroup of
\(\underline{\operatorname{Aut}}(\underline{\mathcal G}^{+})\)
in scan 36; replacing that automorphism construction by the
underlying group \(\underline{\mathfrak S}_M\) would change the
question.  We leave this comparison explicitly at that point.

For a precise first finite deck-quotient test, choose the free
group \(F_2=\langle x,y\rangle\) and the surjection
\[
 q_2:F_2\longrightarrow(\mathbb Z/2)^2,
 \qquad x\longmapsto(1,0),\qquad y\longmapsto(0,1).
\]
More generally let \(q_n\) be the same abelianization map
modulo \(n\), and let \(\mathcal R_{q_n}(F_2)\) denote its
rational relative completion over the constant finite group
\(((\mathbb Z/n)^2)_{\mathbb Q}\), in the universal-property
sense of E.19.  Neither a comparison of
\(\mathcal R_{q_2}(F_2)\), nor a comparison of its quotient
obtained by abelianizing the unipotent kernel, with a specified
version of \(M_{\rho,\sigma,\mathbb Q}\) or
\(N_{\rho,\sigma,\mathbb Q}\) is proved here.
The comparison morphism, compatibility with the chosen group
map and deck quotient, and its image and kernel remain to be
constructed.  This is the first nontrivial \(n\) in this
deck-quotient family; it is distinct from the signed
\(\mathbb F_2\)-base question above.

For the target \(\mathfrak S_3\times\widehat{\mathbb Z}^{\,*}\),
the corresponding group-scheme question must specify the
actual tower fundamental group and its continuous quotient
map, and either the proconstant target over \(\mathbb Q\)
or the compatible system of constant finite targets
\(\mathfrak S_3\times(\mathbb Z/m)^{*}\).
No such tower-source comparison is constructed here, and
the map \(q_n\) is not asserted to provide it.
The proved modular comparison, the finite-stage cocycle
computations, and E.24--E.25 are the partial results for this
group-scheme follow-up; no unrestricted identification of
Fermat-tower completions is claimed.

\subsection{The covering test and the two rational theories}

\begin{statement}{E.21}{Proved}{The Stallings lead, with the precise completions}
Let \(F=\langle x,y\rangle\) be the discrete free group, use
\([x,y]=xyx^{-1}y^{-1}\), and define
\(f(x)=x\), \(f(y)=y[x,y]\).  For every finite \(c\), \(f\)
induces an automorphism of \(F/\gamma_{c+1}F\).  It induces
continuous automorphisms of every pro-\(p\) completion and of
the rational prounipotent completion.  Nevertheless the action
\(x\mapsto(12),y\mapsto(13)\) on three letters is transitive,
while its pullback along \(f\) has orbits \(\{1,2\},\{3\}\).
\end{statement}
\begin{proof}
The induced map on \(F^{\mathrm{ab}}\) is the identity.
The group \(\gamma_iF/\gamma_{i+1}F\) is generated by iterated
commutators of weight \(i\) in \(x,y\), and replacing either
generator by itself times a commutator changes such an expression
only in weight at least \(i+1\).  Thus \(f\) is the identity on
each successive quotient.  Induction in the central extensions
\[
 1\to\gamma_cF/\gamma_{c+1}F\to F/\gamma_{c+1}F
                   \to F/\gamma_cF\to1
\]
proves that it is an automorphism at every stage.  The inverses
are compatible; passing to rational unipotent versions of these
nilpotent quotients and then taking limits gives the asserted
prounipotent automorphism.  This proves the needed special case
of the Stallings conclusion directly.

For a pro-\(p\) group generated by \(x,y\), their images under
\(f\) generate its Frattini quotient, since they agree there
with \(x,y\).  They topologically generate the group: otherwise
their closed subgroup would be contained in an open subgroup
of index \(p\), contradicting that assertion.  Thus the induced
endomorphism is surjective.  A finitely generated profinite
group is Hopfian.  Indeed it has only finitely many open normal
subgroups of any fixed index; inverse image under a surjective
endomorphism is an injective self-map of each such finite set
and hence a permutation.  Its kernel must therefore lie in
every open normal subgroup, and is trivial.  This proves the
pro-\(p\) assertion for every prime.

For \(a=(12)\), \(b=(13)\), direct permutation multiplication
gives \(b[a,b]=a\).  Hence the pulled-back action has both
generators acting by \((12)\).  The program computes the two
orbit decompositions exactly.  This also shows why neither
the individual pro-\(p\) quotients nor the inverse system of
nilpotent quotients can recover this particular finite action:
\(\mathfrak S_3\) is not nilpotent.  No assertion is made here
that every discrete feature of \(F\) is lost, or that this was
Grothendieck's historical reason for introducing \S50.
\end{proof}

\begin{statement}{E.22}{Proved}{The modular envelope detects the three-letter covering}
The quotient \(\mathfrak S\to\mathfrak S_3\) induced by
\(\operatorname{SL}_2(\mathbb Z)\to\operatorname{SL}_2(\mathbb F_2)\)
extends to both \(\underline{\mathfrak S}_M\) and
\(E_{\mathfrak S}\) as a morphism to the constant finite
\(\mathbb Z\)-group scheme \((\mathfrak S_3)_{\mathbb Z}\).
Consequently these envelopes retain the associated three-sheeted
covering.  The unipotent completion of \(\mathfrak S\) does not.
\end{statement}
\begin{proof}
Both \(\lambda_0\) and \(\lambda_1\) are the identity modulo 2,
so the reduction map kills \(\Pi_0\) and factors through
\(H=\mathfrak S/M\).  Its image is all of
\(\operatorname{SL}_2(\mathbb F_2)\simeq\mathfrak S_3\): the
reductions of \(\rho\) and \(\sigma\) have orders 3 and 2 and
generate the six matrices.  Compose the component projection
of the finite-stage model, or the proconstant projection of
the envelope, with this quotient.  The three-letter permutation
representation therefore factors through either scheme.  E.18
shows that the unipotent completion of \(\mathfrak S\) is trivial,
so this nontrivial action cannot factor through it.

For the free group in E.21, its free metabelian quotient in the
modern sense also retains the action because
\(\mathfrak S_3''=1\).  This does not identify that free
metabelian quotient with an object to which \S50 applies: it is
not virtually abelian.  Indeed it surjects onto the two-generated
metabelian group \(\mathbb Z\wr\mathbb Z\).  In that group, if
an abelian finite-index subgroup existed, it would contain a
nonzero power of the shift and a nonzero power of the lamp at
the origin; these do not commute, since the shift moves that
lamp to a different position.  A quotient of a virtually abelian
group is virtually abelian, giving the contradiction.
The manuscript's terminology and the
reviewer's modern use of ``metabelian'' have different scopes.
\end{proof}

\begin{statement}{E.23}{Known}{What the two rational theories retain on this example}
For the pointed classifying space \(B\mathfrak S\), fibrewise
rationalization retains the discrete fundamental group
\(\mathfrak S\) and has no higher homotopy groups.  To\"en's
schematization has fundamental group
\(P_{\mathbb Q}(\mathfrak S)\).  In this example it is
\(BP_{\mathbb Q}(\mathfrak S)\).  The envelope of \S50 is instead
the integral group scheme \(E_{\mathfrak S}\), whose rational
representations are exactly the quasi-unipotent rational
representations of \(\mathfrak S\).  The full pro-algebraic
completion is strictly larger, even for this particular group.
\end{statement}
\begin{proof}
\citet[Theorem~3; repeated as Theorem~16]{E_Rivera2021} identify fibrewise rational
equivalences by the unchanged discrete \(\pi_1\) and the
rationalized higher homotopy groups.  Applying that statement
to the Eilenberg--Mac Lane space \(B\mathfrak S\), whose higher
homotopy groups vanish, gives the first assertion with no
nilpotence assumption.  The same conclusion follows directly
from its loop algebra \(\mathbb Q[\mathfrak S]\) and its
category of all derived local systems; here arbitrary, not
only finite-dimensional, modules are allowed.

\citet[Corollary~3.3.7]{E_Toen2006} identify the fundamental
group of schematization with the full pro-algebraic completion.
For completeness we check that no higher schematic homotopy
is needed for this example.  A finitely generated free abelian
group \(M\) is algebraically good over \(\mathbb Q\): on each
finite-dimensional representation, its group cohomology is
computed by the Koszul complex of the commuting operators
\(g_i-1\).  After extending scalars to an algebraic closure,
decompose into simultaneous generalized characters.  If some
character value is not 1, the corresponding operator is
invertible and contracts the complex.  On the unipotent block,
\(g_i-1\) differs from \(\log g_i\) by a commuting invertible
factor, so the complex is the Lie-algebra cohomology complex
of the additive vector group.  These are exactly the
cohomology complexes for the pro-algebraic completion of
\(M\).  Explicitly, over \(\overline{\mathbb Q}\) this completion
is \(\mathbb G_a^r\times D((\overline{\mathbb Q}^{\,*})^r)\):
simultaneous multiplicative Jordan decomposition gives the
diagonalizable factor, and logarithms of the commuting
unipotent parts give the additive factor.  Its diagonalizable
factor has exact invariants, by the direct grading decomposition,
so precisely the trivial-character block remains in cohomology.
Every finite collection of algebraic characters descends to a
number field and then, by restriction of scalars, to a rational
representation; hence this description applies to the base
change of the rational completion.  Equality of the cohomology
comparisons descends by faithful scalar extension.

The finite extension \(1\to M\to\mathfrak S\to H\to1\)
gives the corresponding extension of pro-algebraic groups.
One can construct it by the pushout of the cocycle of E.12
along \(M\to P_{\mathbb Q}(M)\); density proves its universal
property for all representations.  In both the discrete and
algebraic Hochschild--Serre spectral sequences, taking
\(H\)-invariants is exact because \(|H|=12\) is invertible in
\(\mathbb Q\).  The cohomology comparison for \(M\) therefore
gives the comparison for \(\mathfrak S\), for every finite
dimensional local system.  The criterion of
\citet[Lemma~3.3.2 and Definition~3.4.1]{E_Toen2006} now gives
\((B\mathfrak S\otimes\mathbb Q)^{\mathrm{sch}}
 \simeq BP_{\mathbb Q}(\mathfrak S)\).

Finally, induce the character
\(M\to\mathbb Q^*\), \(\lambda_0\mapsto2\),
\(\lambda_1\mapsto1\), from the index-12 subgroup \(M\)
to \(\mathfrak S\).  This gives a 12-dimensional rational
representation.  The eigenvalues of \(\lambda_0\) include 2,
so it is not quasi-unipotent and cannot factor through the
envelope by E.14.  The matrices, the defining group relations,
and all twelve eigenvalues are computed in
\texttt{E\_envelope.py}.  This proves strictness on the
requested example, rather than importing the unrelated
counterexample used in the earlier GALOIS package.

None of these comparisons says that the envelope kills a
nontrivial element of \(\mathfrak S\).  That group is residually
finite: an element outside \(M\) is detected in \(H\), and a
nonzero element of \(M\) is detected in some \(M/nM\).
Thus its canonical map to the envelope, and consequently to
the full pro-algebraic completion, is injective on the
discrete group.  What differs is the universal representation
problem and the resulting target, not an elementwise kernel
on this particular input.
\end{proof}

\subsection{The automorphism parameters in the first Fermat example}

\begin{statement}{E.24}{Proved}{The fine parameter scheme and the rational automorphism group}
Put \(e_0=(1,0)^t\), \(e_1=(0,1)^t\),
\(A_0=I-R\), and \(B_0=I-S\), with the integral action matrices
of E.15.  In \S51.2, fols.~667--668, C4 scans 39--41, fix
the reading supported by (21), (28), and (29):
\[
 \nu=2\nu'+\nu_2,\qquad
 A_0\alpha-\nu_3e_1=B_0\beta+\nu e_1,
 \qquad \mu=2\nu+1\in\mathbb G_m,
 \quad \nu_2,\nu_3\in\{0,1\}.
                                                        \tag{E.24.1}
\]
For each of the four sign pairs, the resulting fine parameter
scheme is the smooth relative three-dimensional scheme
\[
 \operatorname{Spec}\mathbb Z[a,b,k,
                   (12k+6\nu_2+4\nu_3+1)^{-1}].          \tag{E.24.2}
\]
It is smooth also at 2 and 3.  This assertion concerns the
parameter scheme with conjugators retained.

For every \(\mathbb Q\)-algebra, the functor of automorphisms
of \(\underline{\mathfrak S}_M\) preserving \(\underline M\),
whose two conjugators lie in \(\underline M\) and whose images
satisfy the loop relation in (17), is represented by
\(\mathbb G_a\rtimes\mathbb G_m\).  Its parameters and action are
\[
 \begin{split}
 u_{t,\nu}(\rho)&=(t,\nu-t)\rho,\qquad
 u_{t,\nu}(\sigma)=(t,-t)\sigma,\qquad \mu=1+2\nu,\\
 u_{t,\nu}|_{\underline M}&=\mu I,\qquad
 (t,\mu)(t',\mu')=(t+\mu t',\mu\mu').
 \end{split}                                             \tag{E.24.3}
\]
This connected rational group differs from the finite-stage
envelope, which has twelve geometric connected components.
\end{statement}
\begin{proof}
Write \(\alpha=(a_0,a_1)\), \(\beta=(b_0,b_1)\), and
\(w=\nu'\).  The two equations in (E.24.1) reduce to
\[
 b_0=a_0+a_1+b_1,\qquad 3a_1-2w=\nu_2+\nu_3=:d.
\]
The complete integral parametrization is
\[
 a_0=a,\quad b_1=b,\quad a_1=d+2k,\quad
 w=d+3k,\quad b_0=a+d+2k+b;
\]
its inverse has \(k=w-a_1\).  The unit to invert is
\(\mu=12k+6\nu_2+4\nu_3+1\), proving (E.24.2).
The coefficient matrix on \((\alpha,\beta)\) has maximal-minor
gcd 3, exactly as the marginal calculation on scan 41 suggests.
After the \(\nu'\) column is included, that gcd is 1.  Thus the
failure of surjectivity of the former linear map at 3 is not a
singularity of the latter parameter scheme.  At 2 and 3 the
unit \(\mu\) is automatically nonzero on every fibre; each fixed
sign component has \(p^3\) points.  For \(p>3\) it has
\(p^2(p-1)\) points.  For \(\mathbb Z/q\), \(q=4,8,9\), each
component has \(q^3\) points.  The program
\texttt{E\_automorphisms.py} verifies these counts by enumerating
the pairs \((a_1,w)\), with the two free coordinates included,
for every requested prime up to 13 and each of these moduli.

We distinguish the other legible sign alternatives.  Replacing
the signs in the two formulas by
\(\nu=2\nu'+s_2\nu_2\) and
\(A_0\alpha+s_3\nu_3e_1=B_0\beta+\nu e_1\),
where \(s_2,s_3\in\{\pm1\}\), changes only the integral constant
in \(3a_1-2\nu'=s_2\nu_2-s_3\nu_3\).
Every such parameter scheme remains an open subscheme of
\(\mathbb A^3_{\mathbb Z}\), but its fibres need not agree:
when \(s_3=+1\) and \(\nu_3=1\), its characteristic-three fibre
is empty because \(\mu\equiv0\).  The matrix
\(\tau=\operatorname{diag}(-1,1)\) gives
\(\tau(\rho)\rho^{-1}=\lambda_1^{-1}\),
\(\tau(\sigma)\sigma^{-1}=\omega\), and acts as \(-I\) on
\(M\).  Its parameters \(\alpha=\beta=0\),
\(\nu_2=\nu_3=1\), \(\nu=-1\) agree with the minus sign
in (28), and give a point also in characteristic three.
This is mathematical evidence for that sign with these fixed
conventions, not a final palaeographical decision.
The signs \(s_2=\pm1\) are related by translating \(\nu'\)
by \(\nu_2\); formula (21) fixes the plus convention.
The alternative middle ``\(=\)'' in (32), if read as a chain of
equalities, is not equivalent to (31): take
\(\alpha=\beta=e_0\), \(\nu=\nu_2=\nu_3=0\).
Then \(A_0\alpha=B_0\beta=(1,-1)^t\ne0\), whereas their
difference is zero.  All four coherent sign alternatives are
computed at \(p=2,3,5\).

Over a \(\mathbb Q\)-algebra, \(A_0\) is invertible and
\(\operatorname{Im}B_0=\{(t,-t)\}\).  Consequently the
conjugacy and loop conditions are exactly
\(\eta=(t,-t)\) and \(\zeta=\eta+\nu e_1=(t,\nu-t)\).
Using the affine realization in E.16, direct multiplication
of the images of \(\rho,\sigma\) gives
\[
 u(\rho)^3=u(\sigma)^2=\omega,
 \qquad u(\lambda_0)=\mu e_0,
 \qquad u(\lambda_1)=\mu e_1.
\]
These equalities include their finite components.  The first
ones give the order relations, and the last two commute.
Thus the presentation in E.15 gives a homomorphism on the
discrete group.  Its restriction to \(M\) extends by the
linear map \(\mu I\); that map commutes with \(R,S\).
The pushout construction E.12 therefore gives a group-scheme
endomorphism.  Conversely every additive polynomial over a
\(\mathbb Q\)-algebra is linear: in the equation
\(f(x+y)=f(x)+f(y)\), the coefficient of \(x^{n-1}y\)
annihilates \(na_n\), so \(a_n=0\) for \(n>1\).
The images of the basis of \(M\) therefore determine its
restriction to \(\underline M\), and the finite component
representatives then determine the whole homomorphism.

Composition on the generator images gives
\(t''=t+\mu t'\) and \(\nu''=\nu+\mu\nu'\), whence
\(\mu''=\mu\mu'\).  The inverse has parameters
\((-t/\mu,-\nu/\mu)\).  This proves all functorial assertions
in (E.24.3), as well as the precise difference from the
underlying finite-stage group scheme.  The symbolic relations
and the composition identity are checked in the same program.
\end{proof}

\begin{statement}{E.25}{Refuted}{Generator images do not determine the integral automorphism functor}
The assertion in \S51.2, fol.~665 verso, C4 scan 36, formula (16),
that an automorphism is determined by its values on
\(\rho,\sigma\) is false for the finite-stage model over all
bases.  For every prime \(p\), let
\(A=\mathbb F_p[\epsilon]/(\epsilon^2)\).  In the integral
cocycle coordinates of E.16, the formula
\[
 (v,h)\longmapsto(T(v),h),\qquad
 T(v)=v+\epsilon\bigl(v^{[p]}-v\bigr)                 \tag{E.25.1}
\]
defines a nonidentity automorphism of
\((\underline{\mathfrak S}_M)_A\) which fixes every element
of the discrete group \(\mathfrak S\), and preserves the
loop subgroup \(\mathbb G_a\lambda_0\).  Here \(v^{[p]}\)
means coordinatewise Frobenius.  In particular it satisfies
the positive-sign existential conjugator conditions with
\(\alpha=\beta=0\), yet is invisible to the two-image
parameters.  The smallest example in this family is the
four-element ring \(\mathbb F_2[\epsilon]/(\epsilon^2)\).
Requiring a scalar action on the entire loop subgroup would
exclude this example, but is stronger than the manuscript's
condition that this subgroup be preserved.

Moreover the proposed \(\varepsilon_3(u)\) in scan 37, fol.~666,
formula (18), need not be determined by \(u\) even over
\(\mathbb Q\).  In the order-24 ambient model, the identity
automorphism has both the conjugator \(\alpha=1\) and the
oppositely oriented conjugator
\[
       \alpha=(-1/3,1/3)\tau.
\]
Consequently retaining the conjugator sign defines additional
parameter data, rather than a character on the unqualified
automorphism subgroup.

A correct replacement over an arbitrary base is to impose
that the automorphism preserve \(\underline M\) and restrict
to an \(A\)-linear map there.  Within that subfunctor the
two generator images do determine the automorphism.
Over \(\mathbb Q\)-algebras linearity is automatic, as in E.24.
\end{statement}
\begin{proof}
In characteristic \(p\), \(T\) is additive.  The polynomial
\(U(x)=x-\epsilon(x^p-x)\) is its inverse: expanding
\(T(U(x))\) and \(U(T(x))\) uses \(\epsilon^2=0\) and
\(\epsilon^p=0\), and gives \(x\) in both cases.  It is
nonidentity because the coefficient of \(x^p\) is the
nonzero element \(\epsilon\).  It commutes with every
matrix over \(\mathbb F_p\), since Frobenius fixes their
entries.  In particular it commutes with every action
matrix \(A_h\), and it fixes every integral cocycle value
\(c(h,k)\), whose coordinates reduce to \(\mathbb F_p\).
Applying it to
\[
 (v,h)(w,k)=(v+A_hw+c(h,k),hk)
\]
therefore proves the group identity for (E.25.1) as a
polynomial identity over \(A\).  Each discrete lattice
coordinate also lies in the prime field and is fixed.
The coordinate line generated by \(e_0\) is preserved by
both \(T\) and \(U\), proving the loop-subgroup assertion.
The program reduces symbolic polynomials modulo
\((p,\epsilon^2)\), checks additivity, both inverse identities,
all twelve action matrices, and all 144 cocycle values,
for \(p=2,3,5\).  These are polynomial checks, so the
counterexample does not rely on equality only on the
finite set of \(A\)-valued points.

For the sign assertion, E.24 gives
\(\tau\rho\tau^{-1}\rho^{-1}=\lambda_1^{-1}\).
If \(\alpha=\alpha_0\tau\), its commutator with \(\rho\)
is therefore \((I-R)\alpha_0-e_1\) in the additive normal
subgroup.  Taking \(\alpha_0=(-1/3,1/3)\) makes this zero.
The orientation of \(\alpha_0\tau\) is negative, whereas
the orientation of the identity conjugator is positive.
Both are allowable conjugators in \(\{1,\tau\}\underline M\)
and give the identical value \(u(\rho)=\rho\).
The exact matrix identity and rational vector computation
are checked in \texttt{E\_automorphisms.py}.

For the replacement, agreement on \(\rho,\sigma\) implies
agreement on their words \(\lambda_0,\lambda_1\).
Two \(A\)-linear maps on \(\underline M\) agreeing on its
basis agree identically.  Each component of the model is
a translate of \(\underline M\) by a word in \(\rho,\sigma\),
so the two morphisms agree on every component.  This proves
the replacement without reducedness or characteristic
restrictions.  It also explains why finite-stage
automorphisms and the full envelope's universal property
must be kept separate: the further transition indexed by
\(pM\) kills the nonconstant correction in (E.25.1), in
agreement with the colimit uniqueness of E.14.
\end{proof}
